\section{Running a job}\label{sec-running}
\progname{} uses both MPI and OpenMP parallel programming. Serial, MPI, OpenMP, and hybrid MPI/OpenMP execution are supported. MPI parallelism is the recommended choice for most calculations.

\begin{enumerate}
\item \textbf{Serial execution.} Run the program with a single process and one thread.
\begin{lst-installation}
diabat diabat.inp
\end{lst-installation}
\item \textbf{MPI execution.} Start the required number of MPI processes using the launcher available on the local system:
\begin{lst-installation}
mpirun -np 4 diabat diabat.inp
\end{lst-installation}
In this example, four MPI processes are used.
On an HPC system, use \texttt{mpiexec} or the launcher required by the job scheduler when appropriate.
\item \textbf{OpenMP execution.} Select the number of threads through \texttt{OMP\_NUM\_THREADS} or the \texttt{-nt} option.
\begin{lst-installation}
export OMP_NUM_THREADS=4
diabat diabat.inp
\end{lst-installation}
Alternatively, the command may be written as
\begin{lst-installation}
diabat diabat.inp -nt 4
\end{lst-installation}
OpenMP is used in selected orbital-integral and cube-density calculations and in portions of the optimization for certain diabatization tasks. It does not parallelize every computational step, so MPI execution is usually more effective.
\item \textbf{Hybrid MPI/OpenMP execution.} MPI processes may each use several OpenMP threads.
\begin{lst-installation}
export OMP_NUM_THREADS=2
mpirun -np 4 diabat diabat.inp
\end{lst-installation}
Hybrid execution is generally unnecessary unless the selected calculation spends substantial time in OpenMP-enabled operations or the computing environment favors this arrangement.
\end{enumerate}

If \texttt{OMP\_NUM\_THREADS} and \texttt{-nt} are both provided, the environment variable takes precedence. The default is one OpenMP thread. The installed program prints the supported execution forms with \texttt{diabat --usage}.

\subsection{Standard output and error output}\label{sec-stdout}
\progname{} writes progress information and calculation results to standard output \texttt{stdout}. Diagnostic and error messages are sent to standard error \texttt{stderr}. Both streams normally appear in the \keyemph{terminal}.

Shell redirection can save them to separate files. For example, the following command writes normal output to \texttt{diabat.log} and error messages to \texttt{diabat.err}.
\begin{lst-installation}
diabat diabat.inp > diabat.log 2> diabat.err
\end{lst-installation}
To combine both streams in one file, use
\begin{lst-installation}
diabat diabat.inp > diabat.log 2>&1
\end{lst-installation}
These redirections can also be used for parallel runs.
Individual jobs may additionally create numerical data and wave-function files in their selected output directories.
